在均匀各向同性线弹性半空间 \(z\ge0\)，自由表面 \(z=0\) 无切向、法向牵引。密度 \(\rho>0\)、Lamé 常数 \(\mu>0,\lambda\ge0\)。平面应变小位移 \(\boldsymbol u=(u_x,u_z)\) 满足
\[
\rho\boldsymbol u_{tt}=(\lambda+\mu)\nabla(\nabla\cdot\boldsymbol u)+\mu\nabla^2\boldsymbol u.
\]
设纵横波速分别为 \(c_p=\sqrt{(\lambda+2\mu)/\rho}\)、\(c_s=\sqrt{\mu/\rho}\)。研究沿 \(x\) 传播、波数 \(k>0\)、角频率 \(\omega>0\)、相速 \(v=\omega/k<c_s\) 且随 \(z\to\infty\) 衰减的 Rayleigh 表面波。
(1) 用 \(\boldsymbol u=(\phi_x-\psi_z,\phi_z+\psi_x)\) 和势函数 \(\phi=Ae^{-pkz}e^{i(kx-\omega t)}\)、\(\psi=Be^{-qkz}e^{i(kx-\omega t)}\)，推导正衰减常数 \(p,q\) 及应力自由面关于 \(A,B\) 的完整齐次方程。
(2) 设 \(\zeta=v^2/c_s^2\)、\(\gamma=c_s^2/c_p^2\)。导出 Rayleigh 方程与消根号后的三次式；严格证明 \(0<\gamma\le1/2\) 时恰有一个 \(0<\zeta<1\) 的传播物理根，并排除平方产生的静态零根。
(3) 求归一化幅度未固定的纵横位移复振型，证明表面轨迹两分量相差四分之一周期；用弹性应力功推导局部能量守恒，并由自伴模态变分关系求深度积分的时间平均能流速度。
(4) 对 \(\lambda=0\) 与 \(\lambda\to\infty\) 分别计算 \(v/c_s\) 至至少六位有效数字，求两种势的穿透长度；解释无截止色散以及 \(v<c_s<c_p\) 与体波辐射条件的关系。